\section{Introduction}

Robots are increasingly capable of replacing routine manual and cognitive tasks. However, there remains a strong demand for human cognitive and motor skills that machines cannot yet replicate. At the same time, this technological shift is occurring alongside a growing labor shortage, partly because many workers are unwilling to take jobs that are physically demanding and are expected to be replaced by automation in the future.

As a result, industries such as agriculture, forestry, and heavy-duty mobile machinery are already facing significant shortages of skilled workers \parencite{machado_autonomous_2021}. Moreover, career expectations are increasingly influenced by perceptions of automation. Research shows that when individuals believe a job is likely to be replaced by artificial intelligence, their willingness to pursue that career decreases significantly \parencite{wang_influence_2026}.

Despite the growing level of automation, current robotic systems still struggle with complex and dynamic environments. Therefore, skilled human operators remain essential. However, acquiring these skills is difficult and time-consuming, as it requires extensive repeated practice to develop muscle memory.

In this context, learning theory provides a useful framework. near transfer refers to applying skills in similar contexts, while far transfer refers to applying them in different contexts \parencite{barnett_when_2002}. For near transfer, repeated practice plays a critical role in building procedural knowledge and muscle memory.

Given these challenges, Extended Reality (XR) has emerged as a promising training approach. Previous studies have explored XR in domains such as medical procedures and aviation \parencite{dutescu_virtual_2025}. Meanwhile, Unmanned Aerial Vehicles (UAVs) are widely used in construction, agriculture, logistics, and disaster management \parencite{otto_optimization_2018}. Among these applications, First Person View (FPV) quadcopter piloting represents a particularly demanding skill that relies heavily on precise motor control and spatial awareness. Accordingly, this study uses FPV quadcopter piloting as a case to investigate XR-based training for muscle memory development.

However, existing drone simulators are primarily designed for entertainment rather than structured learning. Consequently, they often present steep learning curves and provide limited support for systematic skill acquisition. In addition, the lack of standardized controller configurations create further barriers \parencite{ebeid_survey_2018}, leading to user frustration even before effective training begins.

Taken together, low career expectations, high preparation effort, and early-stage frustration significantly reduce learner motivation. As a result, many beginners discontinue training before acquiring sufficient muscle memory.

To address these issues, the goal of this study is to design an XR-based training system that lowers the learning barrier for FPV quadcopter piloting. In particular, the proposed system aims to support muscle memory development through AI-assisted co-training and improved interaction design.

Accordingly, this study addresses the following research questions:

\begin{itemize}
\item[] (1) How can FPV quadcopter flight physics be simulated?
\item[] (2) How can a human–AI co-training framework be designed to assist users?
\item[] (3) What interaction methods enable high-fidelity XR prototyping?
\end{itemize}

This study uses a design-based research approach. It applies agile development practices and iterative prototyping to build an XR FPV quadcopter simulator, including physics simulation, AI co-training (PID control and reinforcement learning), and interaction design.